% !PRINTED none
% Part-title / section-opening page: no running head / page number visible
% in any of the three bands (top band is essentially blank). Content:
% the "Achter Abschnitt." part title with its descriptive subtitle, a
% dividing rule, then the "§ 78." section heading and the opening of its
% body text.
\begin{center}
\gesperrt{Achter Abschnitt.}

\gesperrt{Die Gruppe der Transformationsgleichungen und die Gleichung 5ten Grades.}
\end{center}

\noindent\rule{0.3\textwidth}{0.4pt}

\begin{center}
§ 78.

\gesperrt{Die Galoissche Gruppe der Transformationsgleichungen für einen Primzahlgrad.}
\end{center}

Ein eingehenderes algebraisches Studium der Transformationsgleichungen erfordert die Kenntnis ihrer \gesperrt{Galoisschen} Gruppe. Da
wir die Transformationsgleichungen aus den Teilungsgleichungen
hergeleitet haben, deren Gruppe uns bekannt ist (§ 63), so können
wir die Gruppe der Transformationsgleichungen gleichfalls bilden.

Die Transformationsgleichungen ergaben sich (§ 65) dadurch,
daß die Wurzeln der Teilungsgleichungen sich in Reihen einteilen ließen, die durch die Vertauschungen der Gruppe der
Teilungsgleichung nicht auseinandergerissen, sondern nur untereinander vertauscht werden.

Jeder dieser Reihen ordnet sich eine bestimmte Wurzel einer
Transformationsgleichung zu, und die \gesperrt{Gruppe der letzteren
besteht daher aus dem Inbegriff der Vertauschungen, die durch die Gruppe der Teilungsgleichung unter den Reihen hervorgerufen werden.}

Ist der Transformationsgrad $n$ \gesperrt{eine ungerade Primzahl $p$},
so gestattet diese Gruppe eine sehr elegante Darstellung, die zu
weiteren Untersuchungen geeignet ist, und wir halten jetzt diese
Voraussetzung fest.

Es wurde im § 65, (3) bereits die notwendige und hinreichende
Bedingung ermittelt, daß zwei Wurzeln der Teilungsgleichung

\weq{1}{x_{\mu,\mu'},\quad x_{\nu,\nu'}}

in dieselbe Reihe $R$ gehören, nämlich die Kongruenz

\weq{2}{\mu\,\nu'-\nu\,\mu' \equiv 0 \pmod{p}.}
